% ==========================================
% 第 1 章：RoboCup 与 HSL
% ==========================================
\section{RoboCup 与 HSL}

RoboCup 于 1997 年举办首届赛事。其长期目标之一，是推动全自主机器人足球发展到能够与人类顶级球队竞技的水平。人形机器人足球是其中一条技术路线，要求机器人在没有人工遥控的情况下完成感知、决策和运动控制。\official

在这项比赛中，「全自主」意味着机器人需要独立完成比赛决策；「人形」则要求它在双足支撑和运动过程中维持平衡。感知、决策与运动控制不仅要分别工作，还必须在有限的计算资源和时间内协同运行。

\subsection{机器人足球考验什么}

机器人足球同时考验视觉、定位、策略和运动控制。机器人需要从相机图像中识别球与场地标记，估计自身位置，再根据比赛状态选择动作，最后通过关节控制完成移动或射门。这个过程会持续重复，任何一环的延迟或错误都可能传递到后续模块。例如，球的位置估计落后于真实位置时，控制器即使准确执行了目标动作，机器人也可能朝球已经离开的位置移动。

团队通常按这条链路划分工作：感知负责估计机器人、球和场地的位置；策略根据这些信息与比赛状态选择行为；运动控制将目标行为转换成可执行的关节动作。第 4 章按相关技术展开，读者可以先阅读与当前分工直接相关的部分。

比赛规则对自主运行和外部通信有明确限制。软件设计时应确认哪些信息可以从场外接收，并保证比赛策略能够在机载计算机上独立运行。具体限制以当届规则为准。

\subsection{相关组别与技术积累}

RoboCup 从来不止一个联赛。早些年同时铺着五六个方向，各自解决一块问题；最近几年，资源明显向人形足球收拢。

\begin{table}[H]
\centering
\small
\begin{tabularx}{\textwidth}{p{2.6cm} p{3.4cm} X}
\toprule
\textbf{组别} & \textbf{现在的状态} & \textbf{留给 HSL 的东西} \\
\midrule
SPL 标准平台组 & 并入人形体系 & 定位、队伍通信、多机战术的开源代码。B-Human 是这一支留下的完整遗产。 \\
\addlinespace
MSL 中型组 & 停办 & 全向底盘的运动规划和高速击球。双足用不上，但工程方法可借。 \\
\addlinespace
仿真组 & 压缩规模 & 便宜的训练与回归测试环境。它能替代真机跑通流程，替代不了真机调参数。 \\
\addlinespace
@Home 服务组 & 独立发展 & 家庭服务场景，与足球无关。本书不涉及。 \\
\bottomrule
\end{tabularx}
\caption{本书写作时的组别格局。调整每年都有，以当届官网公告为准。}
\end{table}

对开发者而言，组别变化带来的直接影响，是可参考的代码来自不同平台和赛制。人形机器人硬件、接触条件和场地光照都与轮式平台或纯仿真环境不同，因此复用旧方案时，需要先区分通用的软件设计与依赖具体硬件的实现。\advice

阅读早期开源项目时，可以优先研究定位、通信和状态机等系统设计，再判断哪些实现适用于当前平台。步态和平衡控制与硬件结构、执行器特性密切相关，通常不能直接移植。

\subsection{一场比赛的全过程}

参赛流程不仅包括场上的对抗，也包括检录、比赛状态切换和赛后处理。软件需要为每个阶段定义明确行为，尤其要保证机器人在非比赛状态下不会执行不应发生的动作。

\begin{table}[H]
\centering
\small
\begin{tabularx}{\textwidth}{p{2.4cm} X X}
\toprule
\textbf{阶段} & \textbf{场上发生什么} & \textbf{你的代码该做什么} \\
\midrule
赛前检录 & 裁判量尺寸、称重、查传感器，确认合规 & 提前核对尺寸与重量，并为必要的调整留出时间 \\
\addlinespace
INITIAL & 机器人在场外待命，等待入场 & 只收包，不动。所有电机保持安全状态 \\
\addlinespace
READY & 机器人自行走到开球站位 & 按规则到达指定位置，并确认定位状态满足启动条件 \\
\addlinespace
SET & 全部就位，静止 & 冻结底盘，准备开球动作 \\
\addlinespace
PLAYING & 自由对抗 & 感知、决策、运控全速跑；只有这个状态在计时 \\
\addlinespace
半场 & 交换场地，短暂休息 & 坐标整体镜像；重新初始化定位 \\
\addlinespace
FINISHED & 比赛结束 & 按规定进入安全状态，等待工作人员处理 \\
\bottomrule
\end{tabularx}
\caption{一场比赛的阶段与代码的对应关系。状态名和裁判盒一致，见 1.6 节。}
\end{table}

在团队赛制下，机器人需要分配站位、抢球和补位等任务，再根据队友状态进行协调。具体人数与换人规则可能随赛季调整，软件应从配置或裁判状态中读取相关信息，而不是把人数写死在逻辑中。\official

\subsection{场地条件对算法的影响}

\begin{table}[H]
\centering
\small
\begin{tabularx}{\textwidth}{p{2.2cm} p{4.2cm} X}
\toprule
\textbf{场地元素} & \textbf{它是什么} & \textbf{对算法意味着什么} \\
\midrule
草皮 & 人工草，表面可能反光 & 检查白平衡和曝光；如使用颜色阈值，应验证其对场地变化的敏感程度 \\
\addlinespace
白色标线 & 边界、中线、中圈、禁区 & 定位的主要特征来源；线宽固定，可用来估距离 \\
\addlinespace
中圈与罚球点 & 圆形和点状标记 & 圆形对称，容易产生镜像歧义，粒子滤波会收敛到错误一侧 \\
\addlinespace
球门 & 两根立柱加横梁 & 射门目标；立柱细，远了就丢，需要专门的检测策略 \\
\addlinespace
球 & 白底黑斑 & 滚动时外观变化大，检测要抗运动模糊 \\
\addlinespace
光照 & 场馆顶灯，也可能有自然光进入 & 光照可能随时间和场地位置变化，视觉参数应便于检查和调整 \\
\bottomrule
\end{tabularx}
\caption{场地元素、感知特征及其对算法的影响。具体尺寸以当届规则为准。}
\end{table}

现场出现感知或定位问题时，可以先根据表中的对应关系缩小排查范围，再用图像、日志和实测结果确认原因。

\subsection{规则里和写代码有关的六条}

以下列出几类会直接影响软件设计的规则要求。这份概述不能替代当届规则原文，赛前应逐项核对。

\begin{enumerate}[leftmargin=*]
    \item \textbf{全自主。} 比赛期间的自主运行、外部通信和人工干预均受规则约束。应确认程序在允许的机载环境中完成决策，并按队内调试规范保存日志。\official
    \item \textbf{尺寸与重量有上限。} 改机械之前先量一遍。检录不合格，整队白来。
    \item \textbf{传感器有限制。} 激光雷达明确禁用；深度相机在强光下噪点多，检录时属于灰色地带。定位只能靠可见光加里程计，第 4.7 节讲怎么做。
    \item \textbf{换人有代价。} 死球时才能申请，回场要等罚时结束。决策层要能处理「少一个人」的局面，也要能从零重新定位。
    \item \textbf{危险动作会被判罚。} 冲撞、推挤以及守门员持球时间等行为可能受到规则限制。策略层应将相关约束纳入行为选择，并在必要时允许机器人停止接近或主动退让。
    \item \textbf{半场交换场地。} 上下半场换边。代码里进攻方向不能写死，要能整体镜像一次。
\end{enumerate}

传感器限制会影响感知方案的选择；半场交换场地则要求定位、进攻方向和行为策略使用一致的坐标变换约定。

\subsection{裁判盒：机器人听谁的话}

裁判通过 GameController（比赛控制程序，队内简称裁判盒或 GC）广播当前比赛状态。机器人接收并解析这些信息，再据此决定是否允许移动以及何时执行比赛策略。

\begin{itemize}[leftmargin=*]
    \item \textbf{它怎么工作：} 主裁判在笔记本上运行 GC，选择比赛状态；GC 通过局域网把状态包广播到 3838 端口；场上所有机器人同时收到，各自解析。\official
    \item \textbf{它不管什么：} 不知道机器人在哪、球在哪。它只负责计时、计分、发状态、记罚时。你的定位、视觉、决策，它一概不参与。
    \item \textbf{它为什么重要：} 如果机器人收不到或错误解析状态包，就可能在不正确的比赛状态下继续运行。因此，集成测试应覆盖丢包、编号不一致和状态切换等情况。
\end{itemize}

\begin{table}[H]
\centering
\small
\begin{tabularx}{\textwidth}{p{2.4cm} X}
\toprule
\textbf{数据包字段} & \textbf{用途} \\
\midrule
队伍编号、球员编号 & 判断这个包属于哪台机器人；编号错了就全盘收不到 \\
\addlinespace
比赛状态 & INITIAL / READY / SET / PLAYING / FINISHED，以及各种罚时子状态 \\
\addlinespace
剩余时间 & 决定要不要抢最后一脚；也用来判断半场是否结束 \\
\addlinespace
罚时信息 & 谁被罚下、还剩多久，决定场上阵型怎么补 \\
\bottomrule
\end{tabularx}
\caption{裁判盒数据包里和决策相关的字段（字段名以官方仓库的定义为准）。}
\end{table}

\noindent\textbf{首次联调。} 建议按以下顺序逐项确认：

\begin{enumerate}[leftmargin=*]
    \item 下载官方 GameController，选对组别和队伍数量。
    \item 设好 Team Number 和 Player ID，和代码配置里逐位核对。
    \item 电脑和机器人放同一网段，关掉防火墙和虚拟网卡。
    \item 抓包确认：能收到广播，字段和时间都对得上。
    \item 走一遍完整的 INITIAL \(\rightarrow\) READY \(\rightarrow\) SET \(\rightarrow\) PLAYING，看机器人状态跳转。
\end{enumerate}

\noindent\textbf{三个常见坑：}

\begin{itemize}[leftmargin=*]
    \item \textbf{队号或球员号不一致。} 这类配置错误可能导致机器人忽略收到的数据包。联调时应先核对官方配置与程序中的编号。
    \item \textbf{防火墙或多网卡。} 网络配置可能阻止广播包到达程序，表现为间歇性丢包或只有部分机器人能接收。排查时应确认实际使用的网卡、网段和防火墙规则。
    \item \textbf{半场交换后方向错误。} 如果坐标变换不一致，机器人可能朝错误的球门移动。应在统一的坐标约定下处理换边，不要在多个模块中分别修改方向常量。
\end{itemize}

\subsection{五个常见误解}

\begin{description}[leftmargin=0pt,style=nextline]
    \item[「机器人是遥控的吧？」] 比赛要求机器人自主运行，人工干预受到规则限制。场边手持的控制器可能用于安全停机或进入特定保护状态，具体功能应按赛事规则和设备说明确认。
    \item[「必须买官方指定的机器人？」] 规则只要求合规：尺寸、重量、传感器都在规则里写着。自己造一台能过检录，完全可以参赛。
    \item[「仿真跑通了，真机应该也行。」] 仿真通过只能说明程序在当前模型和参数下表现正常。地面摩擦、执行器延迟、结构形变和光照变化都可能造成仿真与真机之间的差异，因此仍需逐步进行实机验证。
    \item[「规则是裁判的事，跟写代码没关系。」] 规则可能影响传感器选型、比赛状态处理、换人策略和坐标变换，因此需要在软件设计阶段一并考虑。
    \item[「代码写得越复杂越强。」] 复杂度本身并不代表性能优势。模块边界清楚、故障可定位、行为可复现，往往更有利于现场调试和长期维护。
\end{description}

\newpage